% Zdroj: PDF priklady-podzim-2022.pdf, fyzická str. 2. Značka: společné okruhy. Zařazeno: M2. Číselné výsledky ověřeny numericky.
\section{Báze a ortogonální báze}
\szzotazka{podzim 2022}{4}{Báze a ortogonální báze}

\noindent\textbf{Zadání.}\par
Uvažujme matici $A = BC$, kde
\[
  B = \begin{pmatrix}
    1 & 3 \\
    2 & 2 \\
    1 & 2 \\
    1 & 0 \\
    1 & -1
  \end{pmatrix},
  \qquad
  C = \begin{pmatrix}
    1 & 2 & 3 & 4 & 5 \\
    5 & 4 & 3 & 2 & 1
  \end{pmatrix}.
\]
\begin{enumerate}
  \item Najděte bázi sloupcového prostoru matice $A$.
  \item Definujte pojem ortogonální báze.
  \item Najděte ortogonální bázi sloupcového prostoru matice $A$.
\end{enumerate}

\medskip
\noindent\textbf{Náčrt řešení.}\par
\begin{enumerate}
  \item Matice $A = BC$ je typu $5 \times 5$, ale $\operatorname{rank}(A) = \operatorname{rank}(B) = 2$. Platí totiž rovnost sloupcových prostorů $\operatorname{col}(A) = \operatorname{col}(B)$: každý sloupec $A$ je tvaru $B \cdot (\text{sloupec } C)$, takže $\operatorname{col}(A) = \{Ax : x \in \mathbb{R}^5\} = \{B(Cx) : x \in \mathbb{R}^5\}$. Protože $C$ má plnou řádkovou hodnost $2$, je $\{Cx : x \in \mathbb{R}^5\} = \mathbb{R}^2$, a tedy $\operatorname{col}(A) = \{By : y \in \mathbb{R}^2\} = \operatorname{col}(B)$. Bázi proto tvoří dva (lineárně nezávislé) sloupce matice $B$:
    \[
      b_1 = (1,2,1,1,1)^{\mathrm{T}}, \qquad b_2 = (3,2,2,0,-1)^{\mathrm{T}}.
    \]
  \item \emph{Ortogonální báze} prostoru (se skalárním součinem) je báze, jejíž každé dva různé vektory jsou na sebe kolmé, tj.\ mají nulový skalární součin.
  \item Použijeme Gramovu--Schmidtovu ortogonalizaci na $b_1, b_2$:
    \[
      u_1 = b_1 = (1,2,1,1,1)^{\mathrm{T}}, \qquad \|u_1\|^2 = 8.
    \]
    Dále $b_2 \cdot u_1 = 8$, takže
    \[
      u_2 = b_2 - \frac{b_2 \cdot u_1}{\|u_1\|^2}\, u_1 = b_2 - \frac{8}{8}\, u_1 = b_2 - u_1 = (2,0,1,-1,-2)^{\mathrm{T}}.
    \]
    Kontrola: $u_1 \cdot u_2 = 0$. Hledaná ortogonální báze je tedy
    \[
      \{(1,2,1,1,1)^{\mathrm{T}},\ (2,0,1,-1,-2)^{\mathrm{T}}\}.
    \]
\end{enumerate}

\medskip
\noindent\textit{(V~PDF bez náčrtu řešení; náčrt doplněn k~výkladu okruhu, výpočet ověřen numericky.)}
